% Chapter 1 — Introduction and mathematical background
%
% Demonstrates: inline math, displayed equations, aligned multi-line
% equations, equation references with \eqref, piecewise (cases),
% matrix–vector notation, acronyms, figures, and citations.

\chapter{Introduction and mathematical background}
\addthumb{Introduction}{\thechapter}{white}{thesisblue}
\clearpage

\section{Motivation and thesis outline}

\lipsum[1]

The performance of a detection algorithm is typically assessed using
the \ac{roc} curve, which plots the true-positive rate against the
false-positive rate as the decision threshold varies~\cite{bishop2006}.
A central design goal in many \ac{ml} pipelines is to maximise the
\ac{snr} of the processed signal so that downstream classifiers
operate on clean features~\cite{shannon1948}.

\lipsum[2]

\section{Inline and displayed equations}

Inline mathematics is enclosed in dollar signs, e.g.\ the energy
$E = mc^{2}$ or the Euler identity $e^{i\pi} + 1 = 0$.
For important standalone results, use a numbered displayed equation:

\begin{equation}
  (f * g)(t) = \int_{-\infty}^{\infty} f(\tau)\, g(t - \tau)\, d\tau.
  \label{eq:ch1-convolution}
\end{equation}

The convolution in~\eqref{eq:ch1-convolution} appears throughout
signal processing and probability theory.
Cross-reference equations with \verb|\eqref|, which automatically adds
parentheses.

\section{Aligned multi-line equations}

Use the \texttt{align} environment when a derivation spans several
lines and you want the equals signs to line up.
Each line receives its own label; suppress unwanted labels with
\verb|\notag|.

\begin{align}
  \hat{f}(\xi) &= \int_{-\infty}^{\infty} f(t)\,
                   e^{-2\pi i \xi t}\, dt,
                   \label{eq:ch1-fourier-fwd} \\
  f(t)         &= \int_{-\infty}^{\infty} \hat{f}(\xi)\,
                   e^{+2\pi i \xi t}\, d\xi.
                   \label{eq:ch1-fourier-inv}
\end{align}

Equations~\eqref{eq:ch1-fourier-fwd}–\eqref{eq:ch1-fourier-inv}
form the Fourier transform pair~\cite{shannon1948}.
A short derivation that needs only one label uses \verb|\notag| on
intermediate lines:

\begin{align}
  \|f + g\|^{2}
    &= \langle f + g,\, f + g \rangle \notag \\
    &= \|f\|^{2} + 2\operatorname{Re}\langle f, g \rangle
       + \|g\|^{2} \notag \\
    &\leq \bigl(\|f\| + \|g\|\bigr)^{2}.
       \label{eq:ch1-triangle}
\end{align}

Inequality~\eqref{eq:ch1-triangle} is the triangle inequality in an
inner-product space.

\section{Piecewise functions and matrix notation}

A piecewise definition uses the \texttt{cases} environment:

\begin{equation}
  \operatorname{ReLU}(x) =
  \begin{cases}
    x & \text{if } x > 0, \\
    0 & \text{otherwise.}
  \end{cases}
  \label{eq:ch1-relu}
\end{equation}

A fully connected layer applies the linear map in~\eqref{eq:ch1-linear}
followed by a non-linearity such as~\eqref{eq:ch1-relu}:

\begin{equation}
  \mathbf{y} = \mathbf{W}\mathbf{x} + \mathbf{b},
  \qquad
  \mathbf{W} \in \mathbb{R}^{m \times n},\;
  \mathbf{x} \in \mathbb{R}^{n},\;
  \mathbf{b} \in \mathbb{R}^{m}.
  \label{eq:ch1-linear}
\end{equation}

\section{Figures}

Every figure must have a caption that fully describes what is shown,
including panel labels (a), (b), axis units, and any abbreviations.
Place the \verb|\label| \emph{inside} or immediately after the
\verb|\caption| so that \verb|\ref| points to the correct number.

\begin{figure}[t]
  \centering
  \includegraphics[width=0.5\linewidth]{example-image}
  \caption{Placeholder figure. Describe what is shown here, including
           any panel labels, axis units, and abbreviations used.
           Adapted from~\protect\cite{bishop2006}.}
  \label{fig:ch1-example}
\end{figure}

An overview of the experimental setup is shown in
Figure~\ref{fig:ch1-example}.
\lipsum[3]

\printbibliography[heading=subbibliography, title={References}]
